% TeXquill feature playground. Start here; see README.md for the walkthrough.
% Open examples/demo as the editor workspace and attach texquill --stdio.
% The unmodified project compiles without intentional errors.
\documentclass[11pt]{article}
\usepackage[T1]{fontenc}
\usepackage{mathtools,amssymb}
\usepackage{xcolor}
\usepackage{etoolbox}
\usepackage{hyperref}

% PATHS: go to definition / open document link on this input or bibliography.
% Rename sections/introduction.tex through an LSP-aware file explorer to try
% workspace/willRenameFiles; the input path should be updated automatically.
% !TeX root = main.tex
% Shared declarations make navigation, signatures and rename cross-file.

% The set of real numbers, for use in math mode.
\newcommand{\R}{\mathbb{R}}

% Format a term with an optional category and a short explanation.
%
% #1: category (defaults to concept).
% #2: term to emphasize.
% #3: explanation printed after the term.
\newcommand{\demoterm}[3][concept]{%
  \textbf{#2} (\emph{#1}): #3\par
}

% Welcome a reader using demoterm; inspect this macro's outgoing calls.
% #1: the reader's name.
\newcommand{\demowelcome}[1]{%
  \demoterm[greeting]{Welcome}{Hello, #1.}%
}

% A panel with an optional heading and a required title.
% #1: heading (defaults to Note); #2: title.
\newenvironment{demopanel}[2][Note]
  {\begin{quote}\textbf{#1: #2}\par}
  {\end{quote}}

\newtheorem{demotheorem}{Theorem}[section]
\DeclarePairedDelimiter{\norm}{\lVert}{\rVert}

% Inspect a literal color here; rename accent to update uses in main.tex.
\definecolor{accent}{HTML}{2563EB}

\newcounter{demostep}
\newlength{\demogap}
\setlength{\demogap}{0.5em}
\newbool{showdetails}
\booltrue{showdetails}


\title{TeXquill LSP Playground}
\author{Example Author}
\date{}

\begin{document}
\maketitle
\tableofcontents

% 1. CROSS-FILE NAVIGATION, COMPLETION, REFERENCES AND RENAME
% On sec:intro or knuth1984 below: hover, definition/declaration, references,
% document highlights, and rename. Preview the edits in the other files.
% For completion, temporarily empty a key or type a prefix inside its braces.
See Section~\ref{sec:intro}, Equation~\eqref{eq:identity}, and
Theorem~\ref{thm:identity}. For background, see~\cite{knuth1984,lamport1994}.

% !TeX root = ../main.tex
\section{Introduction}
\label{sec:intro}

% Navigate here from main.tex, or rename this label and preview all edits.
This section lives in another file. Its label is referenced from the root
document, and its own references point back to Equation~\eqref{eq:identity}
and the editing exercises in Section~\ref{sec:editing}.

% The definition is in a third file: hover, signature help, references, rename.
\demoterm{Language server}{One project connects several source files.}

\begin{demopanel}{Cross-file symbols}
  The declarations are in a separate file.
  Their names and arguments support the same editing features here as in
  the root document. For an introduction to LaTeX, see~\cite{lamport1994}.
\end{demopanel}

% Use this printed sentence as the inverse-search target in a PDF viewer.
This sentence is the source/PDF search target in the introduction subfile.


\section{Mathematics and user declarations}
\label{sec:math}

% 2. HOVER, SIGNATURE HELP AND CALL HIERARCHY
% Hover demoterm to read its documentation from definitions.tex.
% Request signature help inside [result] or either required argument.
% Prepare call hierarchy on demowelcome: outgoing calls lead to demoterm;
% incoming calls on demoterm include demowelcome and direct document calls.
% Rename demoterm to update its declaration and calls in all source files.
\demowelcome{reader}
\demoterm[result]{Identity}{Adding zero preserves a real number.}

\begin{equation}
  x + 0 = x, \qquad x \in \R
  \label{eq:identity}
\end{equation}

% User math commands also support navigation, references and rename.
% Try command completion with \alp, \fra, \demo or \norm in math mode.
\[
  \norm{x} = \sqrt{x^2}, \qquad
  \norm*{\frac{x}{2}} = \frac{\norm{x}}{2}.
\]

% 3. ENVIRONMENT SIGNATURES, LINKED EDITING AND SYMBOL RENAME
% Request signature help after \begin{demopanel}[ or inside its title.
% Edit the name in a begin/end pair with linked editing enabled in the client.
% Rename demopanel to update the declaration and every use across files.
\begin{demopanel}[Example]{A custom environment}
  The panel has an optional heading and a required title.
  Navigate back to Section~\ref{sec:intro} or forward to Section~\ref{sec:editing}.
\end{demopanel}

\begin{demotheorem}[Additive identity]
  \label{thm:identity}
  For every real $x$, we have $x + 0 = x$.
\end{demotheorem}

% 4. COUNTERS, LENGTHS, BOOLEANS AND NAMED COLORS
% Navigate, find references, or rename demostep, demogap, showdetails, accent.
\setcounter{demostep}{0}
\refstepcounter{demostep}
Step \arabic{demostep}: \textcolor{accent}{inspect a named color}.
\par\vspace{\demogap}
\ifbool{showdetails}{This sentence is controlled by a named boolean.}{}

\newpage
\section{Editing and presentation}
\label{sec:editing}

% 5. COLOR PREVIEWS AND PRESENTATIONS
% Open the client color picker on the literal below or on definecolor in
% definitions.tex. A named color reference supports rename and completion.
\textcolor[HTML]{2563EB}{A literal blue color.}
\colorbox{accent!15!white}{A named color mixed with white.}

% 6. DOCUMENT LINKS AND PDF DESTINATIONS
% Open the URL link. Navigate/rename demo-target independently of labels.
\href{https://www.latex-project.org/}{The LaTeX Project}
maintains the document preparation system.
\hypertarget{demo-target}{This paragraph is a named PDF destination.}
Use \hyperlink{demo-target}{this internal link} to return here.
\pdfbookmark[1]{Editing playground}{demo-bookmark}

\subsection{Selection and structural refactoring}
\label{sec:refactor}

% 7. Select the words "reusable phrase" below and request code actions:
% wrap in emph / textbf / quote, or extract a new command in this root file.
% Use expand selection repeatedly inside a ref or a nested command argument.
This reusable phrase is a convenient selection for structural refactoring.

% 8. FORMATTING, FOLDING AND SEMANTIC HIGHLIGHTING
% Remove indentation in this nested list, then format the selection/document.
% Press Enter or type } to try on-type indentation (if enabled by the client).
% Fold a section, environment, comment block, display formula or BibTeX entry.
\begin{itemize}
  \item Commands, environments, labels and citations have semantic tokens.
  \item Nested structures provide selection and folding ranges.
    \begin{enumerate}
      \item Follow \ref{sec:math} to another section.
      \item Look up \cite{shannon1948} by key, author or title in completion.
    \end{enumerate}
\end{itemize}

% Literal examples should remain untouched by symbol rename and diagnostics.
\begin{verbatim}
\demoterm{literal}{not a macro call}
\ref{not-a-real-label} and \cite{not-a-real-key}
\begin{not-a-real-environment}
\end{verbatim}

\subsection{Diagnostics and quick fixes}

% 9. Uncomment ONE exercise at a time, request code actions, then undo.
% Undefined references/citations have static spelling fixes without a build:
% See Section~\ref{sec:intr} and the book~\cite{knuth198}.
%
% A duplicate label reports both the conflict and its related definition:
% \label{sec:intro}
%
% An unclosed group offers a syntax repair:
% \textbf{This group is deliberately unfinished.
%
% A mismatched environment offers a repair; try source.fixAll as well:
% \begin{quote}
%   This environment deliberately has the wrong closing name.
% \end{itemize}
%
% Build-log exercise: uncomment, save and build; inspect the compiler error
% and command spelling fix, then undo and rebuild to clear the diagnostic.
% \textbff{A misspelled command for a compiler diagnostic.}

% ChkTeX 46: enable diagnostics.chktex with arguments ["-w46"], uncomment
% the next line and save. Request a code action to use \( ... \) delimiters.
% An inline formula $x + 1$ can use explicit LaTeX math delimiters instead.

\newpage
\section{Build, bibliography and source/PDF search}
\label{sec:build}

% 10. Build using the code lens or latex.build. Rebuild after source edits.
% Once compilation succeeds, enable the client's inlay hints to see label
% numbers/pages and numeric citations next to ref/eqref/cite arguments.
The introduction is on page~\pageref{sec:intro}; the identity is
Equation~\eqref{eq:identity}. The bibliography includes~\cite{knuth1984,shannon1948}.

% Forward search from this paragraph should land on page 3 or later.
% In the PDF, inverse search on the introduction should open the subfile.
This paragraph is a stable target for forward source-to-PDF search.
Compare the source location with the PDF position, then return to the
introduction using inverse search in your editor's viewer integration.

% 11. Open references.bib: try entry/field/string completion, hover,
% definition and rename on the string abbreviation, citation key rename,
% document/range formatting, outline, folding, semantic tokens and links.
\bibliographystyle{plain}
\bibliography{references}
\end{document}
